\definecolor{promptbg}{gray}{0.95}
\definecolor{promptframe}{gray}{0.5}
\definecolor{afterbg}{HTML}{F0FAF0}
\definecolor{afterframe}{HTML}{2E7D32}
\definecolor{beforebg}{HTML}{FFF5F5}
\definecolor{beforeframe}{HTML}{C62828}
\definecolor{thinkbg}{HTML}{FAFAFA}
\definecolor{thinkframe}{gray}{0.75}
% --- mdframed style definitions (add to preamble) ---
\mdfdefinestyle{promptbox}{%
  backgroundcolor=promptbg,
  linecolor=promptframe,
  linewidth=0.8pt,
  innertopmargin=8pt,
  innerbottommargin=8pt,
  innerleftmargin=8pt,
  innerrightmargin=8pt,
  skipabove=6pt,
  skipbelow=4pt
}
\mdfdefinestyle{afterbox}{%
  backgroundcolor=afterbg,
  linecolor=afterframe,
  linewidth=1pt,
  innertopmargin=8pt,
  innerbottommargin=8pt,
  innerleftmargin=8pt,
  innerrightmargin=8pt,
  skipabove=6pt,
  skipbelow=4pt
}
\mdfdefinestyle{beforebox}{%
  backgroundcolor=beforebg,
  linecolor=beforeframe,
  linewidth=1pt,
  innertopmargin=8pt,
  innerbottommargin=8pt,
  innerleftmargin=8pt,
  innerrightmargin=8pt,
  skipabove=6pt,
  skipbelow=4pt
}
\mdfdefinestyle{thinkbox}{%
  backgroundcolor=thinkbg,
  linecolor=thinkframe,
  linewidth=0.5pt,
  innertopmargin=6pt,
  innerbottommargin=6pt,
  innerleftmargin=6pt,
  innerrightmargin=6pt,
  skipabove=4pt,
  skipbelow=4pt
}
% ============================================================
% Place the following in the body of the document
% ============================================================


\begin{figure*}[t]
\scriptsize
% === Full-width PROMPT ===
\begin{mdframed}[style=promptbox]
\noindent\textbf{Prompt}
You perform a high-throughput experiment on white lupine to find genes contributing to resistance to the fungal disease anthracnose. You receive three candidate genes of unknown function~-- G1, G2, and G3. You create three knock-out mutants (g1, g2, g3) and double-mutants (g1g2, g1g3, g2g3). You know that at least one of these genes is a transcription factor acting upstream of other gene(s). After testing with the pathogen, you receive the following results (100\% = wild-type resistance):

\vspace{2pt}
\begin{tabular}{@{}lr@{\hspace{6pt}}lr@{\hspace{6pt}}lr@{}}
\toprule
Genotype & Resistance & Genotype & Resistance & Genotype & Resistance \\
\midrule
g1 & 75\% & g1g3 & 10\% & g1g2 & 0\% \\
g2 & 0\%  & g2g3 & 0\%  &      &     \\
g3 & 50\% &      &      &      &     \\
\bottomrule
\end{tabular}
\vspace{2pt}

\noindent Which conclusion can you draw?

(A)~G2 is a TF, G1\&G3 show pleiotropy, G1 epistatic towards G3

(B)~G2 is a TF, G1\&G3 show gene redundancy, G1 epistatic towards G3

(C)~G2 is a TF, G1\&G3 share a promoter, G3 epistatic towards G1

(D)~G1 is a TF, G2\&G3 show pleiotropy, G2 epistatic towards G1
\hfill\textbf{\textit{Correct answer:}~(B)}
\end{mdframed}
\vspace{2pt}
% === Two columns: Before \sdrlfootnotesize | After \sdrlfootnotesize ===
\noindent
\begin{minipage}[t]{0.49\textwidth}
% --- Before \sdrlfootnotesize ---
\begin{mdframed}[style=beforebox]
\noindent\textbf{Before \sdrlfootnotesize} \hfill {\colorbox{beforebg}{\textcolor{beforeframe}{\textbf{$\boldsymbol\times$~Incorrect (empty)}}}}

\noindent\textit{Thinking (truncated):}
\begin{mdframed}[style=thinkbox]
\begin{itemize}[leftmargin=12pt, labelsep=4pt, itemsep=0pt, parsep=0pt, topsep=0pt]
\item Wild-type (control): 100\% resistance.
  \begin{itemize}[leftmargin=12pt, labelsep=4pt, itemsep=0pt, parsep=0pt, topsep=0pt]
  \item g1 mutant: 75\% resistance (slight decrease).
  \item g2 mutant: 0\% resistance (complete loss).
  \item g3 mutant: 50\% resistance (moderate decrease).
  \item g1g3 double mutant: 10\% resistance (significant decrease, more than g1 or g3 alone).
  \item g2g3 double mutant: 0\% resistance (same as g2).
  \item g1g2 double mutant: 0\% resistance (same as g2). [...]
  \item If G2 is a transcription factor (TF) that activates G1 and G3:
    \begin{itemize}[leftmargin=12pt, labelsep=4pt, itemsep=0pt, parsep=0pt, topsep=0pt]
    \item If G2 is G, G1 and G3 are targets.
    \item If G2 is knocked out, G1 and G3 are not expressed.
    \item So, g2 would have the same phenotype as g2g1 and g2g3.
    \item This matches the data: g2 (0\%), g2g1 (0\%), g2g3 (0\%).
    \end{itemize}
  \item Now, let's look at G1 and G3 again:
    \begin{itemize}[leftmargin=12pt, labelsep=4pt, itemsep=0pt, parsep=0pt, topsep=0pt]
    \item g1: 75\%
    \item g3: 50\%
    \item g1g3: 10\%
    \item If G1 and G3 act in parallel pathways ( pathways,, loss111111111111111111111 G11 G G11 G G G1 G G1 G G G G G G G1 G1 G G1 G G G G G G G1 G G G G G G111 G111 G G G1111 G111111111 G1111 G G1111111111111111111
    \end{itemize}
  \end{itemize}
\end{itemize}
{\itshape [\ldots repeated 1 tokens until the max tokens budget.]}
\end{mdframed}
\vspace{2pt}
\noindent\textit{Response:}
\begin{mdframed}[style=thinkbox]
\textit{[empty --- no response produced]}
\end{mdframed}
\end{mdframed}
\end{minipage}%
\hfill
\begin{minipage}[t]{0.49\textwidth}
% --- After \sdrlfootnotesize ---
\begin{mdframed}[style=afterbox]
\noindent\textbf{After \sdrlfootnotesize} \hfill {\colorbox{afterbg}{\textcolor{afterframe}{\textbf{\checkmark~Correct}}}}

\noindent\textit{Thinking (truncated):}
\begin{mdframed}[style=thinkbox]
\begin{itemize}[leftmargin=12pt, labelsep=4pt, itemsep=0pt, parsep=0pt, topsep=0pt]
\item Wild-type (control): 100\% resistance.
  \begin{itemize}[leftmargin=12pt, labelsep=4pt, itemsep=0pt, parsep=0pt, topsep=0pt]
  \item g1: 75\% resistance (slightly reduced).
  \item g2: 0\% resistance (completely sensitive).
  \item g3: 50\% resistance (moderately reduced).
  \item g1g3: 10\% resistance (highly sensitive).
  \item g2g3: 0\% resistance (completely sensitive).
  \item g1g2: 0\% resistance (completely sensitive).
  \end{itemize}
\end{itemize}
g2 has 0\% resistance. This means G2 is essential for the resistance. g1g3 has 10\%, much lower than both g1 (75\%) and g3 (50\%). This suggests that G1 and G3 work in parallel pathways (redundancy).

\noindent If G2 is a TF upstream of G1 and G3, then knocking out G2 (g2) would prevent the expression of both G1 and G3. This would explain why g2, g2g3, and g1g2 all have 0\% resistance.

\vspace{2pt}
\noindent\textit{[...re-examines epistasis definitions, checks all options, confirms (B)...]}
\vspace{2pt}

\noindent Given the options, (B) is the only one that correctly identifies G2 as a transcription factor and G1/G3 as redundant.
\end{mdframed}
\vspace{2pt}
\noindent\textit{Response:}
\begin{mdframed}[style=thinkbox]
To determine the correct conclusion regarding the genes G1, G2, and G3, we analyze the experimental resistance data provided:
\begin{enumerate}[leftmargin=12pt, labelsep=4pt, itemsep=0pt, parsep=0pt, topsep=2pt]
\item \textbf{Wild-type (control):} 100\% resistance.
\item \textbf{g1 mutant:} 75\% resistance (G1 contributes slightly to resistance). \textit{[...]}
\end{enumerate}

\noindent\textbf{Analysis of the interactions:}
\begin{itemize}[leftmargin=12pt, labelsep=4pt, itemsep=0pt, parsep=0pt, topsep=2pt]
\item \textbf{Transcription Factor:} If G2 is a transcription factor upstream of both G1 and G3, knocking out G2 would prevent the expression of both G1 and G3, \textit{[...]}
\end{itemize}

\vspace{2pt}
\noindent\textbf{Final Answer: (B)}
\end{mdframed}
\end{mdframed}
\end{minipage}
\caption{\textbf{Example from GPQA-Diamond (Biology) showing model behavior with the DiffusionGemma sampler before and after \sdrl training}. The SFT model begins with \emph{correct} reasoning---identifying G2 as an essential transcription factor and recognizing G1/G3 redundancy---but degenerates into a repetitive token loop (the digit `1'), producing an empty final response. After \sdrl, the model completes the analysis cleanly and outputs the correct answer.}
\label{fig:gpqa-rl-example-biology}
\end{figure*}

\begin{figure*}[t]
\scriptsize
% === Full-width PROMPT ===
\begin{mdframed}[style=promptbox]
\noindent\textbf{Prompt}
\vspace{2pt}\noindent
A quantum mechanical particle of mass $m$ moves in two dimensions in the following potential, as a function of $(r, \theta)$:
%
\[
V(r, \theta) = \tfrac{1}{2}\,k\,r^2 + \tfrac{3}{2}\,k\,r^2\cos^2(\theta)
\]
%
Find the energy spectrum.

\vspace{2pt}
\noindent
\begin{tabular}{@{}ll@{}}
(A) & $E = (n_x + 3n_y + \tfrac{3}{2})\,\hbar\sqrt{k/m}$ \\[2pt]
(B) & $E = (2n_x + n_y + \tfrac{3}{2})\,\hbar\sqrt{k/m}$ \\[2pt]
(C) & $E = (3n_x + 2n_y + \tfrac{1}{2})\,\hbar\sqrt{k/m}$ \\[2pt]
(D) & $E = (2n_x + 3n_y + \tfrac{1}{2})\,\hbar\sqrt{k/m}$
\end{tabular}

\hfill\textbf{\textit{Correct answer:} (B)}
\end{mdframed}
\vspace{2pt}
% === Two columns: Before \sdrlfootnotesize | After \sdrlfootnotesize ===
\noindent
\begin{minipage}[t]{0.49\textwidth}
% --- Before \sdrlfootnotesize ---
\begin{mdframed}[style=beforebox]
\noindent\textbf{Before \sdrlfootnotesize} \hfill {\colorbox{beforebg}{\textcolor{beforeframe}{\textbf{$\boldsymbol\times$~Incorrect (empty)}}}}

\noindent\textit{Thinking (truncated):}
\begin{mdframed}[style=thinkbox]
The potential is given in polar coordinates $(r, \theta)$ as:
$V(r, \theta) = \frac{1}{2}kr^2 + \frac{3}{2}kr^2\cos^2(\theta)$.
We know:

\quad $x = r\cos(\theta)$,

\quad $y = r\sin(\theta)$,

\quad $r^2 = x^2 + y^2$,

\quad $\cos^2(\theta) = x^2/r^2$.

\smallskip
Substitute these into the potential expression:

\quad $V(x,y) = \frac{1}{2}k(x^2+y^2) + \frac{3}{2}kx^2 = 2kx^2 + \frac{1}{2}ky^2$

\textit{[...]}

Plugging in the values for $\omega_x$ and $\omega_y$:

\textit{[...]}

The calculated energy spectrum $E = (2n_x + n_y + \tfrac{3}{2})\,\hbar\sqrt{k/m}$, which matches option (B).

\medskip
\textit{[\ldots re-verifies answer three additional times, all correctly arriving at (B) \ldots]}

\medskip
\noindent Let's re-read the options and
\quad .\qquad = ( = =xx\_nn\_\_33/ \qquad ))))
\smallskip
)))))))222\_\_\_\_y +//)) ))) ))))))))))))))2 2 2222\\
)))))))))))))))))))))))))))))))))))))))))))))))))\\
{\itshape [\ldots repeated closing-parenthesis until max token budget.]}
\end{mdframed}
\vspace{2pt}
\noindent\textit{Response:}
\begin{mdframed}[style=thinkbox]
\textit{[empty --- no response produced]}
\end{mdframed}
\end{mdframed}
\end{minipage}%
\hfill
\begin{minipage}[t]{0.49\textwidth}
% --- After \sdrlfootnotesize ---
\begin{mdframed}[style=afterbox]
\noindent\textbf{After \sdrlfootnotesize} \hfill {\colorbox{afterbg}{\textcolor{afterframe}{\textbf{\checkmark~Correct}}}}

\noindent\textit{Thinking (truncated):}
\begin{mdframed}[style=thinkbox]
The potential is $V(r, \theta) = \frac{1}{2}kr^2 + \frac{3}{2}kr^2\cos^2(\theta)$.

\quad The coordinates are polar $(r, \theta)$.

\quad $x = r\cos\theta$

\quad $y = r\sin\theta$

\quad $r^2 = x^2 + y^2$

\quad $\cos\theta = \frac{x}{r}$

\quad $V(r, \theta) = \frac{1}{2}k(x^2+y^2) + \frac{3}{2}k(x^2+y^2)(\frac{x}{r})^2$

\textit{[\ldots]}

The Hamiltonian is separable into two one-dimensional harmonic oscillators:

\quad $H = H_x + H_y$

\textit{[\ldots]}

For a 1D harmonic oscillator $H = -\frac{\hbar^2}{2m}\frac{d^2}{dx^2} + \frac{1}{2}m\omega^2x^2$, the energy eigenvalues are:
    $E = (n+\frac{1}{2})\hbar\omega$, where $n = 0, 1, 2, ...$

\textit{[\ldots]}

The calculated $E = (2n_x+n_y+3/2)\hbar\sqrt{k/m}$ matches option (B).
   
\textit{[\ldots verifies calculation twice, confirms match with option (B).]}
\end{mdframed}
\smallskip
\noindent\textit{Response (excerpt):}
\begin{mdframed}[style=thinkbox]
To find the energy spectrum of the quantum mechanical particle of mass $m$ in the given potential $V(r, \theta) = \frac{1}{2}kr^2 + \frac{3}{2}kr^2\cos^2(\theta)$, we first convert the potential into Cartesian coordinates $(x, y)$.

1.  \textbf{Convert the potential to Cartesian coordinates:}
    The relationship between polar and Cartesian coordinates is:

\textit{[\ldots]}

    % Substituting these into the potential formula:
    % $V(r, \theta) = \frac{1}{2}k(x^2+y^2) + \frac{3}{2}k(x^2+y^2)\left(\frac{x}{r}\right)^2$

% \textit{[\ldots]}

2.  \textbf{Identify the Hamiltonian:}
    The Hamiltonian for a particle in two dimensions is: \textit{[\ldots]}

    % This is the sum of two independent one-dimensional harmonic oscillators, $H = H_x + H_y$.

3.  \textbf{Calculate the energy eigenvalues for each oscillator:}
    The energy levels of a 1D harmonic oscillator $\frac{1}{2}m\omega^2x^2$ are $E = (n+\frac{1}{2})\hbar\omega$.

\textit{[\ldots]}

4.  \textbf{Find the total energy:}

\textit{[\ldots]}

\quad $E = (2n_x + n_y + \frac{3}{2})\hbar\sqrt{\frac{k}{m}}$

\textbf{Final Answer: (B)}
\end{mdframed}
\end{mdframed}
\end{minipage}
\caption{\textbf{Example from GPQA-Diamond (2D harmonic oscillator energy spectrum) showing model behavior with the DiffusionGemma sampler before and after \sdrl training}. The SFT model begins with a \emph{correct} derivation---converting to Cartesian coordinates, identifying the separable Hamiltonian, and computing the correct energy spectrum---but degenerates into a repetitive token loop (closing parentheses), producing an empty final response. After \sdrl, the model completes the same derivation cleanly and outputs the correct answer.}
\label{fig:gpqa-rl-example-quantum}
\end{figure*}
